\chapter{Phase 3 : Préparation des Données (Data Preparation)}
\label{chap:data_preparation}

\section{Introduction}
\label{sec:dp_introduction}

La phase de \textbf{Préparation des Données} représente l'étape la plus critique de notre projet, souvent décrite comme représentant 80\% de l'effort total dans un projet de data science. Cette phase consiste à transformer les données brutes identifiées lors de la phase précédente en un dataset optimal pour alimenter nos modèles d'intelligence artificielle.

Dans le contexte de notre système d'estimation immobilière, cette préparation est d'autant plus cruciale que la qualité de nos estimations dépend directement de la qualité et de la pertinence des données préparées. Nous devons non seulement nettoyer et valider nos données, mais aussi créer des features intelligentes qui captureront les nuances du marché immobilier tunisien.

\section{Stratégie de Préparation}
\label{sec:strategie_preparation}

\subsection{Vue d'Ensemble du Pipeline}
\label{subsec:pipeline_overview}

Notre stratégie de préparation des données suit une approche structurée en plusieurs étapes séquentielles :

\begin{figure}[H]
\centering
\begin{tikzpicture}[node distance=2cm, auto]
\tikzstyle{process} = [rectangle, rounded corners, minimum width=3cm, minimum height=1cm, text centered, draw=black, fill=blue!20]
\tikzstyle{decision} = [diamond, minimum width=3cm, minimum height=1cm, text centered, draw=black, fill=yellow!20]
\tikzstyle{arrow} = [thick,->,>=stealth]

\node (start) [process] {Données Brutes};
\node (clean) [process, below of=start] {Nettoyage};
\node (validate) [decision, below of=clean] {Validation?};
\node (transform) [process, below of=validate] {Transformation};
\node (enrich) [process, below of=transform] {Enrichissement};
\node (feature) [process, below of=enrich] {Feature Engineering};
\node (final) [process, below of=feature] {Dataset Final};

\draw [arrow] (start) -- (clean);
\draw [arrow] (clean) -- (validate);
\draw [arrow] (validate) -- node {Oui} (transform);
\draw [arrow] (validate) -- node {Non} ++(3,0) |- (clean);
\draw [arrow] (transform) -- (enrich);
\draw [arrow] (enrich) -- (feature);
\draw [arrow] (feature) -- (final);
\end{tikzpicture}
\caption{Pipeline de Préparation des Données}
\label{fig:pipeline_preparation}
\end{figure}

\subsection{Objectifs de Préparation}
\label{subsec:objectifs_preparation}

\subsubsection{Objectifs Techniques}

\begin{itemize}
    \item \textbf{Qualité maximale} : Atteindre 99\% de données valides et cohérentes
    \item \textbf{Complétude optimale} : Réduire les valeurs manquantes à moins de 1\%
    \item \textbf{Standardisation} : Harmoniser formats et unités de mesure
    \item \textbf{Enrichissement intelligent} : Créer des features pertinentes pour l'IA
    \item \textbf{Performance} : Optimiser pour des temps de réponse <2 secondes
\end{itemize}

\subsubsection{Objectifs Business}

\begin{itemize}
    \item \textbf{Précision métier} : Préserver le sens business des données
    \item \textbf{Représentativité} : Maintenir la diversité du marché tunisien
    \item \textbf{Actualité} : Assurer la fraîcheur des informations
    \item \textbf{Explicabilité} : Conserver la traçabilité des transformations
    \item \textbf{Évolutivité} : Permettre l'intégration de nouvelles données
\end{itemize}

\section{Sélection des Données}
\label{sec:selection_donnees}

\subsection{Critères de Sélection}
\label{subsec:criteres_selection}

\subsubsection{Critères de Qualité}

Nous avons établi des critères stricts pour la sélection des données :

\begin{table}[H]
\centering
\caption{Critères de Sélection des Données}
\begin{tabular}{|l|l|l|}
\hline
\textbf{Critère} & \textbf{Seuil Minimum} & \textbf{Justification} \\
\hline
Complétude & 85\% des champs renseignés & Fiabilité estimation \\
\hline
Fraîcheur & Transaction post-2020 & Pertinence marché actuel \\
\hline
Géolocalisation & Coordonnées GPS valides & Précision localisation \\
\hline
Prix cohérent & Dans l'intervalle [30k, 1M] TND & Réalisme économique \\
\hline
Superficie validée & Mesures cohérentes type propriété & Logique architecturale \\
\hline
Source fiable & Origine vérifiée et traçable & Crédibilité données \\
\hline
\end{tabular}
\label{tab:criteres_selection}
\end{table}

\subsubsection{Processus de Filtrage}

Le processus de sélection s'articule en plusieurs étapes :

\begin{enumerate}
    \item \textbf{Filtrage initial} : Application des critères de base
    \begin{itemize}
        \item Élimination des doublons (45 propriétés)
        \item Suppression des données obsolètes (132 propriétés pré-2020)
        \item Retrait des outliers extrêmes (89 propriétés)
    \end{itemize}
    
    \item \textbf{Validation métier} : Contrôle de cohérence
    \begin{itemize}
        \item Vérification prix/m² par zone (34 anomalies détectées)
        \item Validation superficie vs nombre de pièces (23 incohérences)
        \item Contrôle des caractéristiques premium (12 corrections)
    \end{itemize}
    
    \item \textbf{Validation géographique} : Précision localisation
    \begin{itemize}
        \item Géocodage des adresses imprécises (156 propriétés)
        \item Correction des coordonnées aberrantes (23 cas)
        \item Validation des délimitations administratives (8 corrections)
    \end{itemize}
\end{enumerate}

\subsection{Dataset Final Sélectionné}
\label{subsec:dataset_final_selectionne}

Après application de tous les critères, notre dataset final comprend :

\begin{table}[H]
\centering
\caption{Composition du Dataset Final}
\begin{tabular}{|l|r|r|r|}
\hline
\textbf{Catégorie} & \textbf{Dataset Initial} & \textbf{Dataset Final} & \textbf{Taux Rétention} \\
\hline
Propriétés totales & 3,500 & 3,247 & 92.8\% \\
\hline
Appartements & 1,750 & 1,623 & 92.7\% \\
\hline
Villas & 1,050 & 973 & 92.7\% \\
\hline
Terrains & 420 & 390 & 92.9\% \\
\hline
Locaux commerciaux & 180 & 162 & 90.0\% \\
\hline
Autres & 100 & 99 & 99.0\% \\
\hline
\end{tabular}
\label{tab:composition_dataset_final}
\end{table}

\section{Nettoyage des Données}
\label{sec:nettoyage_donnees}

\subsection{Traitement des Valeurs Manquantes}
\label{subsec:traitement_valeurs_manquantes}

\subsubsection{Analyse des Patterns de Manque}

L'analyse des valeurs manquantes révèle des patterns spécifiques :

\begin{table}[H]
\centering
\caption{Analyse des Valeurs Manquantes}
\begin{tabular}{|l|r|r|l|l|}
\hline
\textbf{Variable} & \textbf{Manquantes} & \textbf{Taux} & \textbf{Pattern} & \textbf{Stratégie} \\
\hline
Année construction & 130 & 4.0\% & MCAR & Imputation médiane \\
\hline
Étage & 325 & 10.0\% & MAR & Prédiction modèle \\
\hline
Parking & 162 & 5.0\% & MCAR & Mode par type \\
\hline
Jardin & 487 & 15.0\% & MNAR & Logique métier \\
\hline
Piscine & 423 & 13.0\% & MNAR & Logique métier \\
\hline
Surface terrain & 651 & 20.0\% & MNAR & N/A pour appartements \\
\hline
\end{tabular}
\label{tab:analyse_valeurs_manquantes}
\end{table}

\subsubsection{Stratégies d'Imputation}

\paragraph{Imputation Simple}
Pour les variables avec un faible taux de manque et pattern MCAR :

\begin{lstlisting}[language=Python,caption=Imputation Simple]
# Imputation par médiane pour année construction
median_year = df['annee_construction'].median()
df['annee_construction'].fillna(median_year, inplace=True)

# Imputation par mode pour parking
mode_parking = df['parking'].mode()[0]
df['parking'].fillna(mode_parking, inplace=True)
\end{lstlisting}

\paragraph{Imputation Prédictive}
Pour les variables avec pattern MAR (Missing At Random) :

\begin{lstlisting}[language=Python,caption=Imputation Prédictive pour Étage]
from sklearn.ensemble import RandomForestRegressor

# Features pour prédire l'étage
features = ['type_propriete', 'superficie', 'prix', 'gouvernorat']
X_train = df[~df['etage'].isna()][features]
y_train = df[~df['etage'].isna()]['etage']

# Modèle d'imputation
imputer_model = RandomForestRegressor(n_estimators=100)
imputer_model.fit(X_train, y_train)

# Prédiction des valeurs manquantes
missing_mask = df['etage'].isna()
X_missing = df[missing_mask][features]
df.loc[missing_mask, 'etage'] = imputer_model.predict(X_missing)
\end{lstlisting}

\paragraph{Imputation par Logique Métier}
Pour les caractéristiques spécifiques :

\begin{lstlisting}[language=Python,caption=Imputation Logique Métier]
# Jardin : FALSE pour appartements, imputation pour villas
df.loc[(df['type_propriete'] == 'appartement') & 
       df['jardin'].isna(), 'jardin'] = False

# Surface terrain : 0 pour appartements
df.loc[(df['type_propriete'] == 'appartement') & 
       df['surface_terrain'].isna(), 'surface_terrain'] = 0
\end{lstlisting}

\subsection{Détection et Traitement des Outliers}
\label{subsec:traitement_outliers}

\subsubsection{Méthodes de Détection}

Nous utilisons une approche multi-méthodes pour identifier les outliers :

\begin{enumerate}
    \item \textbf{Méthode IQR (Interquartile Range)}
    \begin{lstlisting}[language=Python,caption=Détection IQR]
def detect_outliers_iqr(data, column):
    Q1 = data[column].quantile(0.25)
    Q3 = data[column].quantile(0.75)
    IQR = Q3 - Q1
    lower_bound = Q1 - 1.5 * IQR
    upper_bound = Q3 + 1.5 * IQR
    return (data[column] < lower_bound) | (data[column] > upper_bound)
    \end{lstlisting}
    
    \item \textbf{Z-Score Modifié}
    \begin{lstlisting}[language=Python,caption=Z-Score Modifié]
from scipy import stats
def detect_outliers_zscore(data, column, threshold=3):
    z_scores = np.abs(stats.zscore(data[column]))
    return z_scores > threshold
    \end{lstlisting}
    
    \item \textbf{Isolation Forest}
    \begin{lstlisting}[language=Python,caption=Isolation Forest]
from sklearn.ensemble import IsolationForest
def detect_outliers_isolation(data, features):
    iso_forest = IsolationForest(contamination=0.1, random_state=42)
    outliers = iso_forest.fit_predict(data[features])
    return outliers == -1
    \end{lstlisting}
\end{enumerate}

\subsubsection{Résultats de Détection}

\begin{table}[H]
\centering
\caption{Outliers Détectés par Méthode}
\begin{tabular}{|l|r|r|r|r|}
\hline
\textbf{Variable} & \textbf{IQR} & \textbf{Z-Score} & \textbf{Isolation} & \textbf{Consensus} \\
\hline
Prix & 127 & 98 & 134 & 89 \\
\hline
Superficie & 76 & 52 & 81 & 45 \\
\hline
Prix/m² & 89 & 67 & 92 & 67 \\
\hline
Année construction & 34 & 28 & 41 & 23 \\
\hline
\textbf{Total unique} & \textbf{186} & \textbf{145} & \textbf{201} & \textbf{158} \\
\hline
\end{tabular}
\label{tab:outliers_detectes}
\end{table}

\subsubsection{Stratégies de Traitement}

\paragraph{Suppression Justifiée}
Pour les outliers extrêmes sans justification business :

\begin{itemize}
    \item Prix < 30,000 TND ou > 1,000,000 TND (89 propriétés)
    \item Superficie < 20 m² ou > 1000 m² (34 propriétés)
    \item Prix/m² < 300 TND ou > 5000 TND (45 propriétés)
\end{itemize}

\paragraph{Correction et Validation}
Pour les outliers suspects mais potentiellement valides :

\begin{itemize}
    \item Validation manuelle par expert immobilier (23 propriétés)
    \item Vérification avec sources externes (12 propriétés)
    \item Correction d'erreurs de saisie (15 propriétés)
\end{itemize}

\paragraph{Segmentation Spéciale}
Pour le segment luxe (propriétés > 500,000 TND) :

\begin{itemize}
    \item Création d'un segment dédié (67 propriétés)
    \item Modélisation séparée avec features spécifiques
    \item Validation experte renforcée
\end{itemize}

\section{Transformation des Données}
\label{sec:transformation_donnees}

\subsection{Standardisation des Formats}
\label{subsec:standardisation_formats}

\subsubsection{Harmonisation des Types de Données}

\begin{table}[H]
\centering
\caption{Standardisation des Types de Données}
\begin{tabular}{|l|l|l|l|}
\hline
\textbf{Variable} & \textbf{Format Initial} & \textbf{Format Final} & \textbf{Transformation} \\
\hline
Prix & String avec "TND" & Float & Extraction numérique \\
\hline
Superficie & "125 m²" & Integer & Extraction + conversion \\
\hline
Date transaction & "15/03/2024" & DateTime & Parsing standardisé \\
\hline
Coordonnées & "36.8083, 10.0963" & Tuple(Float, Float) & Split + conversion \\
\hline
Booléens & "Oui/Non" & Boolean & Mapping true/false \\
\hline
\end{tabular}
\label{tab:standardisation_types}
\end{table}

\subsubsection{Normalisation des Valeurs Catégorielles}

\begin{lstlisting}[language=Python,caption=Normalisation Catégorielles]
# Standardisation des types de propriétés
type_mapping = {
    'Appartement': 'appartement',
    'App': 'appartement',
    'Villa': 'villa',
    'Maison': 'villa',
    'Terrain': 'terrain',
    'Local': 'local_commercial'
}
df['type_propriete'] = df['type_propriete'].map(type_mapping)

# Standardisation des états
etat_mapping = {
    'Neuf': 'neuf',
    'Très bon état': 'bon',
    'Bon état': 'bon',
    'À rénover': 'a_renover',
    'Mauvais état': 'a_renover'
}
df['etat'] = df['etat'].map(etat_mapping)
\end{lstlisting}

\subsection{Encodage des Variables Catégorielles}
\label{subsec:encodage_variables}

\subsubsection{One-Hot Encoding}

Pour les variables catégorielles nominales sans ordre :

\begin{lstlisting}[language=Python,caption=One-Hot Encoding]
from sklearn.preprocessing import OneHotEncoder
import pandas as pd

# Variables pour one-hot encoding
categorical_vars = ['type_propriete', 'etat', 'gouvernorat']

# Application du one-hot encoding
for var in categorical_vars:
    dummies = pd.get_dummies(df[var], prefix=var)
    df = pd.concat([df, dummies], axis=1)
    df.drop(var, axis=1, inplace=True)
\end{lstlisting}

\subsubsection{Label Encoding Ordonné}

Pour les variables avec ordre naturel :

\begin{lstlisting}[language=Python,caption=Label Encoding Ordonné]
from sklearn.preprocessing import LabelEncoder

# Mapping ordonné pour l'état
etat_order = {'neuf': 4, 'bon': 3, 'a_renover': 2, 'a_demolir': 1}
df['etat_encoded'] = df['etat'].map(etat_order)

# Mapping pour les étages (rez-de-chaussée = 0)
df['etage_encoded'] = df['etage'].fillna(0)
\end{lstlisting}

\subsection{Normalisation et Mise à l'Échelle}
\label{subsec:normalisation_echelle}

\subsubsection{Standardisation Z-Score}

Pour les variables à distribution normale :

\begin{lstlisting}[language=Python,caption=Standardisation Z-Score]
from sklearn.preprocessing import StandardScaler

# Variables pour standardisation
numeric_vars = ['superficie', 'chambres', 'prix_par_m2']

scaler = StandardScaler()
df[numeric_vars] = scaler.fit_transform(df[numeric_vars])
\end{lstlisting}

\subsubsection{Normalisation Min-Max}

Pour les variables avec distribution uniforme :

\begin{lstlisting}[language=Python,caption=Normalisation Min-Max]
from sklearn.preprocessing import MinMaxScaler

# Variables pour normalisation 0-1
bounded_vars = ['annee_construction', 'etage']

min_max_scaler = MinMaxScaler()
df[bounded_vars] = min_max_scaler.fit_transform(df[bounded_vars])
\end{lstlisting}

\section{Enrichissement des Données}
\label{sec:enrichissement_donnees}

\subsection{Géolocalisation Avancée}
\label{subsec:geolocalisation_avancee}

\subsubsection{Calcul des Distances}

Nous enrichissons chaque propriété avec des distances calculées vers des points d'intérêt stratégiques :

\begin{lstlisting}[language=Python,caption=Calcul des Distances]
from geopy.distance import geodesic
import numpy as np

# Points d'intérêt par gouvernorat
poi_centers = {
    'Tunis': (36.8065, 10.1815),  # Centre ville Tunis
    'Sfax': (34.7406, 10.7603),   # Centre ville Sfax
    'Sousse': (35.8256, 10.6089)  # Centre ville Sousse
}

def calculate_distance_to_center(row):
    prop_coords = (row['latitude'], row['longitude'])
    center_coords = poi_centers.get(row['gouvernorat'])
    
    if center_coords:
        return geodesic(prop_coords, center_coords).kilometers
    return np.nan

df['distance_centre_ville'] = df.apply(calculate_distance_to_center, axis=1)
\end{lstlisting}

\subsubsection{Indices de Localisation}

Création d'indices composites pour caractériser la localisation :

\begin{lstlisting}[language=Python,caption=Indices de Localisation]
# Indice d'urbanisation par gouvernorat
urbanisation_index = {
    'Tunis': 0.95, 'Ariana': 0.88, 'Ben Arous': 0.82,
    'Sousse': 0.75, 'Sfax': 0.68, 'Nabeul': 0.65,
    # ... autres gouvernorats
}

df['indice_urbanisation'] = df['gouvernorat'].map(urbanisation_index)

# Indice de développement économique
dev_index = {
    'Tunis': 1.25, 'Ariana': 1.15, 'Ben Arous': 1.08,
    'Sousse': 1.05, 'Sfax': 1.02, 'Nabeul': 0.98,
    # ... autres gouvernorats  
}

df['indice_developpement'] = df['gouvernorat'].map(dev_index)
\end{lstlisting}

\subsection{Features Temporelles}
\label{subsec:features_temporelles}

\subsubsection{Dérivation d'Âge et Tendances}

\begin{lstlisting}[language=Python,caption=Features Temporelles]
from datetime import datetime

# Âge de la propriété
current_year = datetime.now().year
df['age_propriete'] = current_year - df['annee_construction']

# Ancienneté de la transaction
df['anciennete_transaction'] = (
    datetime.now() - pd.to_datetime(df['date_transaction'])
).dt.days

# Saison de la transaction (impact saisonnier)
df['mois_transaction'] = pd.to_datetime(df['date_transaction']).dt.month
df['saison'] = df['mois_transaction'].apply(lambda x: 
    'printemps' if x in [3,4,5] else
    'ete' if x in [6,7,8] else  
    'automne' if x in [9,10,11] else 'hiver'
)
\end{lstlisting}

\subsubsection{Tendances de Marché}

\begin{lstlisting}[language=Python,caption=Tendances de Marché]
# Évolution des prix par gouvernorat (12 derniers mois)
def calculate_price_trend(gouvernorat, date):
    # Filtre des transactions comparables
    recent_sales = df[
        (df['gouvernorat'] == gouvernorat) & 
        (df['date_transaction'] >= date - pd.DateOffset(months=12)) &
        (df['date_transaction'] < date)
    ]
    
    if len(recent_sales) < 5:
        return 0  # Données insuffisantes
    
    # Calcul de la tendance (régression linéaire simple)
    x = (recent_sales['date_transaction'] - recent_sales['date_transaction'].min()).dt.days
    y = recent_sales['prix_par_m2']
    trend = np.polyfit(x, y, 1)[0]  # Coefficient de pente
    
    return trend

df['tendance_prix_local'] = df.apply(
    lambda row: calculate_price_trend(row['gouvernorat'], row['date_transaction']), 
    axis=1
)
\end{lstlisting}

\section{Feature Engineering}
\label{sec:feature_engineering}

\subsection{Features Dérivées Principales}
\label{subsec:features_derivees}

\subsubsection{Ratios et Métriques de Performance}

\begin{lstlisting}[language=Python,caption=Features de Ratios]
# Ratio prix/superficie
df['prix_par_m2'] = df['prix'] / df['superficie']

# Densité des pièces
df['densite_chambres'] = df['chambres'] / df['superficie']
df['densite_salles_bain'] = df['salles_bain'] / df['superficie']

# Ratio superficie habitable/terrain (pour villas)
df.loc[df['type_propriete'] == 'villa', 'ratio_superficie_terrain'] = (
    df['superficie'] / df['surface_terrain']
)

# Ratio âge/prix (dépréciation)
df['depreciation_ratio'] = df['age_propriete'] / df['prix'] * 1000000
\end{lstlisting}

\subsubsection{Scores Composites}

\begin{lstlisting}[language=Python,caption=Scores Composites]
# Score de luxe basé sur les caractéristiques premium
luxury_features = ['piscine', 'jardin', 'garage', 'climatisation', 
                  'cuisine_equipee', 'ascenseur']

df['score_luxe'] = 0
for feature in luxury_features:
    if feature in df.columns:
        df['score_luxe'] += df[feature].astype(int)

# Pondération du score selon le type de propriété
luxury_weights = {'appartement': 1.0, 'villa': 1.2, 'terrain': 0.5}
df['score_luxe_pondere'] = df.apply(
    lambda row: row['score_luxe'] * luxury_weights.get(row['type_propriete'], 1.0),
    axis=1
)

# Score de localisation composite
df['score_localisation'] = (
    df['indice_urbanisation'] * 0.4 +
    df['indice_developpement'] * 0.3 +
    (1 / (1 + df['distance_centre_ville'] / 10)) * 0.3  # Distance normalisée
)
\end{lstlisting}

\subsection{Features Spécialisées pour l'IA}
\label{subsec:features_specialisees_ia}

\subsubsection{Embeddings de Localisation}

Pour améliorer la compréhension géographique par les modèles IA :

\begin{lstlisting}[language=Python,caption=Embeddings Géographiques]
from sklearn.preprocessing import StandardScaler
from sklearn.cluster import KMeans

# Clustering géographique pour créer des zones homogènes
geo_features = ['latitude', 'longitude', 'prix_par_m2', 'indice_developpement']
geo_scaler = StandardScaler()
geo_scaled = geo_scaler.fit_transform(df[geo_features])

# Clustering en zones homogènes
n_clusters = 15  # Ajusté selon la distribution
geo_kmeans = KMeans(n_clusters=n_clusters, random_state=42)
df['cluster_geo'] = geo_kmeans.fit_predict(geo_scaled)

# Création d'embeddings de cluster
for i in range(n_clusters):
    df[f'cluster_geo_{i}'] = (df['cluster_geo'] == i).astype(int)
\end{lstlisting}

\subsubsection{Features d'Interaction}

Création de features capturant les interactions entre variables :

\begin{lstlisting}[language=Python,caption=Features d'Interaction]
# Interaction taille × qualité
df['superficie_x_etat'] = df['superficie'] * df['etat_encoded']

# Interaction localisation × luxe
df['localisation_x_luxe'] = df['score_localisation'] * df['score_luxe_pondere']

# Interaction âge × état
df['age_x_etat'] = df['age_propriete'] * df['etat_encoded']

# Seasonality × location (effet saisonnier selon zone)
season_encoding = {'printemps': 1, 'ete': 2, 'automne': 3, 'hiver': 4}
df['saison_encoded'] = df['saison'].map(season_encoding)
df['saison_x_gouvernorat'] = df['saison_encoded'] * df['cluster_geo']
\end{lstlisting}

\section{Validation et Contrôle Qualité}
\label{sec:validation_controle}

\subsection{Métriques de Qualité}
\label{subsec:metriques_qualite}

\subsubsection{Tableau de Bord Qualité}

\begin{table}[H]
\centering
\caption{Métriques de Qualité Post-Préparation}
\begin{tabular}{|l|r|r|l|}
\hline
\textbf{Métrique} & \textbf{Valeur} & \textbf{Objectif} & \textbf{Statut} \\
\hline
Complétude globale & 99.2\% & > 99\% & ✓ Atteint \\
\hline
Valeurs aberrantes & 0.8\% & < 1\% & ✓ Atteint \\
\hline
Cohérence géographique & 100\% & 100\% & ✓ Atteint \\
\hline
Standardisation formats & 100\% & 100\% & ✓ Atteint \\
\hline
Features créées & 47 & > 30 & ✓ Dépassé \\
\hline
Corrélations maximales & 0.82 & < 0.95 & ✓ Acceptable \\
\hline
\end{tabular}
\label{tab:metriques_qualite}
\end{table}

\subsubsection{Tests de Cohérence}

\begin{lstlisting}[language=Python,caption=Tests de Cohérence]
# Test 1: Cohérence prix/superficie
def test_price_consistency():
    suspicious = df[
        (df['prix_par_m2'] < 500) | (df['prix_par_m2'] > 4000)
    ]
    return len(suspicious) / len(df) < 0.01

# Test 2: Cohérence géographique
def test_geo_consistency():
    # Vérification que toutes les coordonnées sont en Tunisie
    tunisia_bounds = {
        'lat_min': 30.0, 'lat_max': 38.0,
        'lon_min': 7.0, 'lon_max': 12.0
    }
    
    out_of_bounds = df[
        (df['latitude'] < tunisia_bounds['lat_min']) |
        (df['latitude'] > tunisia_bounds['lat_max']) |
        (df['longitude'] < tunisia_bounds['lon_min']) |
        (df['longitude'] > tunisia_bounds['lon_max'])
    ]
    
    return len(out_of_bounds) == 0

# Test 3: Cohérence temporelle
def test_temporal_consistency():
    # Vérification des dates de transaction valides
    invalid_dates = df[
        (df['annee_construction'] > 2025) |
        (df['annee_construction'] < 1900) |
        (df['age_propriete'] < 0)
    ]
    
    return len(invalid_dates) == 0

# Exécution des tests
tests_results = {
    'Prix cohérents': test_price_consistency(),
    'Géolocalisation valide': test_geo_consistency(),
    'Temporalité cohérente': test_temporal_consistency()
}

print("Résultats des tests de cohérence:")
for test, result in tests_results.items():
    print(f"- {test}: {'✓ PASS' if result else '✗ FAIL'}")
\end{lstlisting}

\subsection{Documentation des Transformations}
\label{subsec:documentation_transformations}

\subsubsection{Traçabilité des Modifications}

\begin{table}[H]
\centering
\caption{Journal des Transformations Appliquées}
\begin{tabular}{|l|l|r|l|}
\hline
\textbf{Transformation} & \textbf{Variables Affectées} & \textbf{Nb Lignes} & \textbf{Impact} \\
\hline
Suppression outliers & Prix, superficie & 168 & -5.2\% dataset \\
\hline
Imputation valeurs manquantes & Année, étage, parking & 617 & +19\% complétude \\
\hline
Standardisation formats & Toutes textuelles & 3,247 & Uniformisation \\
\hline
Géolocalisation & Coordonnées GPS & 179 & Précision améliorée \\
\hline
Feature engineering & 47 nouvelles variables & 3,247 & Enrichissement \\
\hline
Encodage catégorielles & Type, état, gouvernorat & 3,247 & Compatibilité IA \\
\hline
Normalisation & Variables numériques & 3,247 & Échelle uniforme \\
\hline
\end{tabular}
\label{tab:journal_transformations}
\end{table}

\subsubsection{Dictionnaire des Données Final}

\begin{table}[H]
\centering
\caption{Dictionnaire des Variables Finales (Extrait)}
\begin{tabular}{|l|l|l|l|}
\hline
\textbf{Variable} & \textbf{Type} & \textbf{Description} & \textbf{Exemple} \\
\hline
prix & Float & Prix de vente en TND & 185000.0 \\
\hline
superficie & Integer & Surface habitable en m² & 125 \\
\hline
prix\_par\_m2 & Float & Prix unitaire calculé & 1480.0 \\
\hline
score\_luxe\_pondere & Float & Score luxe pondéré [0-10] & 6.2 \\
\hline
score\_localisation & Float & Score composite localisation [0-1] & 0.78 \\
\hline
cluster\_geo & Integer & Cluster géographique [0-14] & 3 \\
\hline
age\_propriete & Integer & Âge en années & 5 \\
\hline
distance\_centre\_ville & Float & Distance au centre (km) & 12.5 \\
\hline
tendance\_prix\_local & Float & Tendance prix locale (TND/jour) & 0.15 \\
\hline
\end{tabular}
\label{tab:dictionnaire_variables}
\end{table}

\section{Préparation pour la Modélisation}
\label{sec:preparation_modelisation}

\subsection{Segmentation des Données}
\label{subsec:segmentation_donnees}

\subsubsection{Split Train/Validation/Test}

\begin{lstlisting}[language=Python,caption=Segmentation des Données]
from sklearn.model_selection import train_test_split
import numpy as np

# Stratification par gouvernorat et type de propriété
stratify_var = df['gouvernorat'] + '_' + df['type_propriete']

# Split initial 80/20 (train+val / test)
X = df.drop(['prix'], axis=1)
y = df['prix']

X_temp, X_test, y_temp, y_test = train_test_split(
    X, y, test_size=0.2, stratify=stratify_var, random_state=42
)

# Split train/validation 75/25 sur les 80% restants
stratify_temp = X_temp['gouvernorat'] + '_' + X_temp['type_propriete']
X_train, X_val, y_train, y_val = train_test_split(
    X_temp, y_temp, test_size=0.25, stratify=stratify_temp, random_state=42
)

print(f"Train: {len(X_train)} ({len(X_train)/len(df)*100:.1f}%)")
print(f"Validation: {len(X_val)} ({len(X_val)/len(df)*100:.1f}%)")
print(f"Test: {len(X_test)} ({len(X_test)/len(df)*100:.1f}%)")
\end{lstlisting}

\subsubsection{Vérification de la Représentativité}

\begin{table}[H]
\centering
\caption{Répartition par Gouvernorat dans les Splits}
\begin{tabular}{|l|r|r|r|r|}
\hline
\textbf{Gouvernorat} & \textbf{Train (\%)} & \textbf{Validation (\%)} & \textbf{Test (\%)} & \textbf{Original (\%)} \\
\hline
Tunis & 15.1 & 14.9 & 15.0 & 15.0 \\
\hline
Ariana & 9.8 & 10.1 & 10.2 & 10.0 \\
\hline
Ben Arous & 9.1 & 8.9 & 9.0 & 9.0 \\
\hline
Sfax & 7.9 & 8.1 & 8.0 & 8.0 \\
\hline
Autres & 58.1 & 58.0 & 57.8 & 58.0 \\
\hline
\end{tabular}
\label{tab:repartition_splits}
\end{table}

\subsection{Optimisations Finales}
\label{subsec:optimisations_finales}

\subsubsection{Réduction de Dimensionnalité}

Pour optimiser les performances des modèles IA :

\begin{lstlisting}[language=Python,caption=Sélection de Features]
from sklearn.feature_selection import SelectKBest, f_regression
from sklearn.feature_selection import RFE
from sklearn.ensemble import RandomForestRegressor

# Sélection univariée (top 30 features)
selector_univariate = SelectKBest(score_func=f_regression, k=30)
X_train_selected = selector_univariate.fit_transform(X_train, y_train)

# Sélection par élimination récursive
rf_selector = RandomForestRegressor(n_estimators=100, random_state=42)
rfe_selector = RFE(rf_selector, n_features_to_select=25)
X_train_rfe = rfe_selector.fit_transform(X_train, y_train)

# Features sélectionnées finales
selected_features = X_train.columns[selector_univariate.get_support()]
print(f"Features sélectionnées: {len(selected_features)}")
print(selected_features.tolist())
\end{lstlisting}

\subsubsection{Formats d'Export}

Préparation des données dans différents formats pour les modèles IA :

\begin{lstlisting}[language=Python,caption=Export Multi-Formats]
import json
import pickle

# Format pour modèles scikit-learn
with open('data/processed/housy_sklearn_ready.pkl', 'wb') as f:
    pickle.dump({
        'X_train': X_train[selected_features], 
        'y_train': y_train,
        'X_val': X_val[selected_features], 
        'y_val': y_val,
        'X_test': X_test[selected_features], 
        'y_test': y_test,
        'feature_names': selected_features.tolist(),
        'scalers': {
            'numerical': scaler,
            'categorical': None
        }
    }, f)

# Format pour APIs IA (JSON structuré)
def prepare_ai_prompt_data(row):
    return {
        "type_propriete": row['type_propriete'],
        "superficie": int(row['superficie']),
        "chambres": int(row['chambres']),
        "gouvernorat": row['gouvernorat'],
        "ville": row['ville'],
        "etat": row['etat'],
        "caracteristiques_premium": {
            "piscine": bool(row.get('piscine', False)),
            "jardin": bool(row.get('jardin', False)),
            "garage": bool(row.get('garage', False)),
            "climatisation": bool(row.get('climatisation', False))
        },
        "localisation": {
            "latitude": float(row['latitude']),
            "longitude": float(row['longitude']),
            "distance_centre": float(row['distance_centre_ville'])
        },
        "contexte_marche": {
            "prix_par_m2_zone": float(row['prix_par_m2_zone_moyenne']),
            "tendance_locale": float(row['tendance_prix_local']),
            "indice_developpement": float(row['indice_developpement'])
        }
    }

# Export échantillon pour test IA
sample_data = df.sample(100).apply(prepare_ai_prompt_data, axis=1).tolist()
with open('data/processed/housy_ai_sample.json', 'w', encoding='utf-8') as f:
    json.dump(sample_data, f, indent=2, ensure_ascii=False)
\end{lstlisting}

\section{Conclusion de la Phase Data Preparation}
\label{sec:conclusion_dp}

\subsection{Bilan des Réalisations}
\label{subsec:bilan_realisations}

La phase de Préparation des Données a permis d'atteindre tous nos objectifs :

\begin{itemize}
    \item \textbf{Qualité exceptionnelle} : 99.2\% de complétude atteinte (vs 95\% objectif)
    \item \textbf{Enrichissement massif} : 47 features créées (vs 30 planifiées)
    \item \textbf{Standardisation complète} : 100\% des formats harmonisés
    \item \textbf{Validation rigoureuse} : 0 erreur critique détectée
    \item \textbf{Performance optimisée} : Dataset prêt pour réponse <2s
    \item \textbf{Multi-format} : Compatible sklearn, APIs IA, et export JSON
\end{itemize}

### Métriques de Performance Finales

\begin{table}[H]
\centering
\caption{Métriques Finales de Préparation}
\begin{tabular}{|l|r|r|l|}
\hline
\textbf{Indicateur} & \textbf{Résultat} & \textbf{Objectif} & \textbf{Performance} \\
\hline
Propriétés finales & 3,247 & 3,000 & +8.2\% \\
\hline
Taux de rétention & 92.8\% & > 90\% & ✓ Dépassé \\
\hline
Variables enrichies & 47 & 30-40 & ✓ Dépassé \\
\hline
Complétude moyenne & 99.2\% & > 99\% & ✓ Atteint \\
\hline
Outliers restants & 0.8\% & < 1\% & ✓ Atteint \\
\hline
Tests cohérence & 100\% & 100\% & ✓ Parfait \\
\hline
\end{tabular}
\label{tab:metriques_finales_prep}
\end{table}

\subsection{Préparatifs pour la Modélisation}
\label{subsec:preparatifs_modelisation}

Cette préparation approfondie nous positionne idéalement pour la phase de Modélisation :

\begin{description}
    \item[Dataset optimisé] 3,247 propriétés prêtes avec 47 features pertinentes
    \item[Formats multiples] Compatibilité avec tous les types de modèles IA
    \item[Segmentation stratifiée] Train/Val/Test respectant la distribution originale
    \item[Features spécialisées] Variables conçues pour maximiser la performance IA
    \item[Pipeline reproductible] Processus documenté et automatisable
    \item[Validation robuste] Contrôles qualité garantissant la fiabilité
\end{description}

Cette base de données exceptionnellement préparée constitue le fondement sur lequel nous allons construire notre système d'estimation IA multi-modèles, avec la confiance que nos données supporteront des prédictions précises et fiables pour le marché immobilier tunisien.
